\documentclass[12pt,a4paper,UTF8]{ctexart}
\usepackage[a4paper,top=22mm,bottom=22mm,left=22mm,right=22mm,headheight=15pt]{geometry}
\newcommand{\Needspace}[1]{\par\begingroup\skip0=\lastskip\vskip-\skip0\dimen0=#1\relax\advance\dimen0 by\skip0\vskip\dimen0\penalty-100\vskip-\dimen0\vskip\skip0\penalty9999\endgroup}
\usepackage{amsmath,amssymb,mathtools,bm}
\usepackage{graphicx,xcolor,booktabs,longtable,array,enumitem,fancyhdr,caption,float,microtype,setspace}
\usepackage[hidelinks,unicode]{hyperref}
\usepackage[most]{tcolorbox}
\definecolor{ContestBlue}{HTML}{17365D}
\hypersetup{pdftitle={第080套},pdfauthor={物理竞赛教学资料}}
\setstretch{1.24}
\setlength{\parindent}{2em}
\setlength{\parskip}{0.35em}
\setlength{\emergencystretch}{2em}
\setlist{itemsep=0.35em,topsep=0.4em}
\captionsetup{font=small,labelfont=bf,justification=centering}
\setcounter{secnumdepth}{0}
\pagestyle{fancy}
\fancyhf{}
\fancyhead[L]{\small 第080套}
\fancyfoot[C]{\thepage}
\renewcommand{\headrulewidth}{0.4pt}
\ctexset{
  section={format=\Large\bfseries\color{ContestBlue},beforeskip=1.2em,afterskip=0.7em},
  subsection={format=\large\bfseries\color{ContestBlue},beforeskip=1.0em,afterskip=0.5em},
  subsubsection={format=\normalsize\bfseries,beforeskip=0.8em,afterskip=0.35em}
}
\usepackage{calc}
\newcounter{none}
\let\originalsubsection\subsection
\renewcommand{\subsection}{\Needspace{6\baselineskip}\originalsubsection}
\makeatletter
\newcommand{\tightlist}{\setlength{\itemsep}{0pt}\setlength{\parskip}{0pt}}
\newcommand{\maxwidth}{\ifdim\Gin@nat@width>\linewidth\linewidth\else\Gin@nat@width\fi}
\newcommand{\maxheight}{\ifdim\Gin@nat@height>0.72\textheight 0.72\textheight\else\Gin@nat@height\fi}
\newcommand{\pandocbounded}[1]{\begingroup\setkeys{Gin}{width=\maxwidth,height=\maxheight,keepaspectratio}#1\endgroup}
\makeatother
\begin{document}
\begin{titlepage}
  \centering
  \vspace*{58mm}
  {\Huge\bfseries\color{ContestBlue} 第080套\par}
  \vfill
  \vspace*{28mm}
\end{titlepage}
\section{第 1 题}\label{ux7b2c-1-ux9898}

空间中有两个固定点电荷 \(q\) 和 \(\alpha q\)，二者相距 \(2r\)。一质量为
\(m\) 的自由空心均匀导体球，其外半径为 \(r\)、空腔半径为
\(r/2\)。初始时，点电荷 \(q\) 位于空腔球心处。不计重力。

\begin{enumerate}
\def\labelenumi{\arabic{enumi}.}
\tightlist
\item
  导体球能够在初始位置保持静止平衡，求 \(\alpha\)。
\item
  在第 1 问条件下，使导体球沿两点电荷连线向右偏移 \(xr\)，其中
  \(x\ll1\)。求导体球所受合力 \(F\) 至 \(x\)
  的最低非零阶；规定向右为正。
\item
  判断初始平衡是否稳定；若稳定，求平衡位置附近的小振动周期。
\end{enumerate}

\begin{figure}[H]
\centering
\includegraphics[width=0.603\linewidth,height=\textheight,keepaspectratio,alt={图 1-1　带球形空腔的导体球及两个固定点电荷的初始位置。}]{../figures/question-01.pdf}
\caption*{图 1-1　带球形空腔的导体球及两个固定点电荷的初始位置。}
\end{figure}

\newpage

\section{第 2 题}\label{ux7b2c-2-ux9898}

将一张长纸条的右端扭转 \(180^\circ\)
后与左端粘合，可形成莫比乌斯带。现制作一个相同拓扑的电阻网络：先由
\(n-1=4\) 个相邻矩形框及两条端边构成平面梯形网络，每条线段的电阻均为
\(R\)；再将 \(A_0\) 与 \(B_5\) 粘合、\(B_0\) 与 \(A_5\) 粘合，得到
\(n=5\) 个网格的莫比乌斯电阻网络。

\begin{enumerate}
\def\labelenumi{\arabic{enumi}.}
\tightlist
\item
  求 \(A_0\)、\(B_0\) 间的等效电阻。
\item
  求 \(A_1\)、\(B_1\) 间的等效电阻。
\item
  对一般很大的 \(n\)，求 \(A_0\)、\(B_0\) 间等效电阻的极限，并比较其与
  \(n=5\) 的结果。
\end{enumerate}

\begin{figure}[H]
\centering
\includegraphics[width=0.9\linewidth,height=\textheight,keepaspectratio,alt={图 2-1　粘合前的五格梯形电阻网络及端点识别关系。}]{../figures/question-02.pdf}
\caption*{图 2-1　粘合前的五格梯形电阻网络及端点识别关系。}
\end{figure}

\newpage

\section{第 3 题}\label{ux7b2c-3-ux9898}

某双星系统中，一颗带磁场的致密星体产生周期约 \(4.84\,\mathrm s\) 的 X
射线脉冲。观测显示脉冲每隔 \(2.1\,\mathrm d\)
隐遁约四分之一周期，说明双星轨道平面近似沿视线方向。可见伴星为红巨星，质量约
\(18.4M_\odot\)、半径约 \(13.08R_\odot\)。设发生掩食时两星中心距可估为
\(a\simeq\sqrt2R_\star\)。取
\(M_\odot\simeq2.0\times10^{30}\,\mathrm{kg}\)，\(R_\odot\simeq7.0\times10^8\,\mathrm m\)。

\begin{enumerate}
\def\labelenumi{\arabic{enumi}.}
\tightlist
\item
  已知中子尺度约为 \(10^{-15}\,\mathrm m\)，估算中子星平均密度。
\item
  根据观测数据估算致密星质量，并由第 1 问密度估算其半径。
\item
  设一个氘核从红巨星表面出发，最终落到中子星表面。忽略广义相对论修正，以经典引力势估算到达表面时的速度。
\end{enumerate}

\begin{figure}[H]
\centering
\includegraphics[width=0.621\linewidth,height=\textheight,keepaspectratio,alt={图 3-1　红巨星与致密星构成的近侧视双星系统及估算轨道尺度。}]{../figures/question-03.pdf}
\caption*{图 3-1　红巨星与致密星构成的近侧视双星系统及估算轨道尺度。}
\end{figure}

\newpage

\section{第 4 题}\label{ux7b2c-4-ux9898}

质量为 \(m\) 的小木块在半径为 \(r\)
的水平圆形轨道内滑动，同时与水平底面和竖直圆壁保持接触。木块尺寸远小于
\(r\)，底面和圆壁与木块之间的动摩擦因数均为 \(\mu\)，木块初速度为
\(v_0\)。

\begin{enumerate}
\def\labelenumi{\arabic{enumi}.}
\tightlist
\item
  当木块已转过角度 \(\theta\)，其动能为 \(E\)。求 \(dE/d\theta\)。
\item
  若木块恰好转过一周后停止，求 \(v_0\)。
\end{enumerate}

\begin{figure}[H]
\centering
\includegraphics[width=0.9\linewidth,height=\textheight,keepaspectratio,alt={图 4-1　木块同时接触水平底面和竖直圆壁并沿圆周滑动。}]{../figures/question-04.pdf}
\caption*{图 4-1　木块同时接触水平底面和竖直圆壁并沿圆周滑动。}
\end{figure}

\newpage

\section{第 5 题}\label{ux7b2c-5-ux9898}

一绝缘介质流体被限制在矩形环形凹槽中均匀流动，泵控制流速；恒定磁场
\(\bm B\) 垂直流体运动平面。两块固定金属板 \(A,B\) 浸入流体，面积均为
\(S\)、相距 \(d\)，并通过外部负载相连。介质相对介电常数为
\(\varepsilon_r\)，忽略边缘效应和回路电流产生的磁场，介质中无自由移动电荷。

\begin{enumerate}
\def\labelenumi{\arabic{enumi}.}
\tightlist
\item
  流体以恒速 \(v_0\)
  运动。空载时求介质两侧的极化电荷面密度和正负号。若外接电阻 \(R\)，且
  \(t=0\) 时两金属板无自由电荷，求 \(U_{AB}(t)\)、\(U_{AB}(\infty)\)
  及电阻消耗总能量 \(W\)。
\item
  系统先在泵持续工作且外电流为零的状态。\(t'=0\)
  时关泵，流速缓慢线性下降：\(v(t')=v_0-at'\)，直至 \(t'_0=v_0/a\)
  后为零。求 \(U_{AB}(t')\) 并作图。
\end{enumerate}

\begin{figure}[H]
\centering
\includegraphics[width=0.9\linewidth,height=\textheight,keepaspectratio,alt={图 5-1　介质流体环路、金属板、磁场方向和外接负载。}]{../figures/question-05.pdf}
\caption*{图 5-1　介质流体环路、金属板、磁场方向和外接负载。}
\end{figure}

\newpage

\section{第 6 题}\label{ux7b2c-6-ux9898}

不透明薄片上打有许多圆孔，孔心按正方形格点排列，晶格常数为
\(d\)。用波长为 \(\lambda\) 的单色激光垂直照射，在薄片后方距离 \(L\)
处放置平行光屏。已知 \(L\gg d\gg\lambda\)。

\begin{enumerate}
\def\labelenumi{\arabic{enumi}.}
\tightlist
\item
  描述光屏上的夫琅禾费衍射图样及主极大位置。
\item
  若薄片沿 \(x\) 轴方向压缩，压缩比例为 \(N\)，衍射图样如何变化？
\end{enumerate}

\begin{figure}[H]
\centering
\includegraphics[width=0.898\linewidth,height=\textheight,keepaspectratio,alt={图 6-1　正方形孔阵及横、纵晶格周期 d。}]{../figures/question-06.pdf}
\caption*{图 6-1　正方形孔阵及横、纵晶格周期 \(d\)。}
\end{figure}

\newpage

\section{第 7 题}\label{ux7b2c-7-ux9898}

真空中有一均匀球形薄膜，内部封有 \(1\,\mathrm{mol}\)
单原子理想气体。薄膜各向同性，单位长度张力恒为 \(\sigma_0\)，面密度为
\(\lambda\)；初始气体温度为 \(T_0\)，系统处于平衡。

\begin{enumerate}
\def\labelenumi{\arabic{enumi}.}
\tightlist
\item
  求初始半径 \(r_0\)。
\item
  薄膜绝热，通过内部电阻丝向系统输入热量 \(Q\)，求达到新平衡后的半径
  \(r_1\)。
\item
  在初态附近给予薄膜整体径向各向同性微扰，判断平衡稳定性；若稳定，求小振动周期。
\end{enumerate}

\begin{figure}[H]
\centering
\includegraphics[width=0.558\linewidth,height=\textheight,keepaspectratio,alt={图 7-1　真空中的球形薄膜及内部一摩尔单原子理想气体。}]{../figures/question-07.pdf}
\caption*{图 7-1　真空中的球形薄膜及内部一摩尔单原子理想气体。}
\end{figure}

\newpage

\section{第 8 题}\label{ux7b2c-8-ux9898}

一静质量为 \(m_0\) 的粒子在地面系中以 \(\beta c=0.8c\)
运动。它突然从外界获得能量而动量不变，裂变成两个静质量均为 \(m_0\)
的粒子；在地面系中，两末态粒子的速度方向相互垂直。求在与初始粒子同速运动的参考系，即初态零动量系中，外界输入能量
\(\Delta E'\) 的允许范围。

\begin{figure}[H]
\centering
\includegraphics[width=0.564\linewidth,height=\textheight,keepaspectratio,alt={图 8-1　地面系中两末态动量相互垂直且矢量和等于初始总动量。}]{../figures/question-08.pdf}
\caption*{图 8-1　地面系中两末态动量相互垂直且矢量和等于初始总动量。}
\end{figure}
\end{document}
